\section*{Chapitre \ref{ch:b3:inorganic-materials} --- Matériaux inorganiques et nanomatériaux}

\begin{solution}{exo:b3:inorganic-materials:1}
$x^2 \propto t$~: une couche trois fois plus épaisse demande neuf fois plus de temps, \qty{90}{h}.
\end{solution}

\begin{solution}{exo:b3:inorganic-materials:2}
Hydrolyse~: $\ce{Si(OC2H5)4 + 4H2O -> Si(OH)4 + 4C2H5OH}$.\\
Condensation~: $\ce{2Si(OH)4 -> Si2O(OH)6 + H2O}$.\\
Bilan~: $\ce{Si(OC2H5)4 + 2H2O -> SiO2 + 4C2H5OH}$.
\end{solution}

\begin{solution}{exo:b3:inorganic-materials:3}
Formateurs~: $\ce{SiO2}$, $\ce{B2O3}$. Modificateurs~: $\ce{Na2O}$, $\ce{CaO}$. Chaque ion oxyde de $\ce{Na2O}$ rompt un pont Si--O--Si
en deux extrémités Si--O$^-$, chacune associée à un ion $\ce{Na+}$~: le réseau est coupé, le liquide fondu
s'écoule plus facilement et fond plus bas.
\end{solution}

\begin{solution}{exo:b3:inorganic-materials:4}
Douze pentagones (comme pour tout \omterm{def:b3:inorganic-materials:carbon}{fullerène}). $V = 70$, $E = 3V/2 = 105$ liaisons, $F = 2 - V + E = 37$
faces, donc $37 - 12 = 25$ hexagones.
\end{solution}

\begin{solution}{exo:b3:inorganic-materials:5}
$\sqrt2\,(r_B + r_O) = 1.4142 \times \qty{200.5}{pm} = \qty{283.5}{pm}$. $\ce{SrTiO3}$~: $284/283.5 = 1.00$, la pérovskite cubique idéale.
$\ce{BaTiO3}$~: $301/283.5 = 1.06$, titane trop petit pour son octaèdre, décentré~: \omterm{def:b3:inorganic-materials:ferroelectric}{ferroélectrique}.
$\ce{CaTiO3}$~: $274/283.5 = 0.97$, calcium trop petit pour sa cavité~: octaèdres basculés, symétrie
abaissée.
\end{solution}

\begin{solution}{exo:b3:inorganic-materials:6}
Si/Al $= 106/86 = 1.23$. $M = 86 \times 23.0 + 86 \times 27.0 + 106 \times 28.1 + 384 \times 16.0 = \qty{13422.6}{g/mol}$~; 43 ions $\ce{Ca^2+}$ par
formule~: $43/\qty{13422.6}{g/mol} = \qty{3.20}{mmol/g}$.
\end{solution}

\begin{solution}{exo:b3:inorganic-materials:7}
Distance entre plus proches voisins $d = a/\sqrt2 = \qty{288.4}{pm}$. Un agrégat de $n$ couches mesure $(2n + 1)d$ de diamètre~:
$2n + 1 \approx 5/0.2884 = 17.3$, donc $n = 8$ (\qty{4.9}{nm}). $N(8) = 2057$ atomes, dont $10 \times 64 + 2 = 642$ en surface~:
31\,\%.
\end{solution}

\begin{solution}{exo:b3:inorganic-materials:8}
$h^2/(8\mu R^2)$ avec $\mu = \qty{9.109e-32}{kg}$~: \qty{0.94}{eV} pour \qty{2.0}{nm} et \qty{0.42}{eV} pour \qty{3.0}{nm}. Émission à
$1.74 + 0.94 = \qty{2.68}{eV}$, \qty{463}{nm} (bleu), et à $1.74 + 0.42 = \qty{2.16}{eV}$, \qty{574}{nm} (jaune)~: la plus petite
boîte émet la lumière la plus bleue.
\end{solution}

\begin{solution}{exo:b3:inorganic-materials:9}
$(10, 10)$~: fauteuil, $n - m = 0$, métallique. $(10, 0)$~: zigzag, 10 n'est pas un multiple de 3, \omterm{def:b3:solid-state:classes}{semi-conducteur}.
$(9, 0)$~: zigzag, métallique. $(12, 8)$~: ni l'un ni l'autre (chiral), $n - m = 4$, \omterm{def:b3:solid-state:classes}{semi-conducteur}.
\end{solution}

\begin{solution}{exo:b3:inorganic-materials:10}
$M = 4 \times 65.4 + 13 \times 16.0 + 24 \times 12.0 + 12 \times 1.0 = \qty{769.6}{g/mol}$~; $a^3 = (\qty{25.82e-8}{cm})^3 = \qty{1.721e-20}{cm^3}$~;
$\rho = 8 \times 769.6/(\num{6.022e23} \times \num{1.721e-20}) = \qty{0.594}{g.cm^{-3}}$. Un terrain mesure \qty{7140}{m^2}~: un gramme en couvre
$3000/7140 = 0.42$, et une cuillerée de quelques grammes plus d'un terrain
entier.
\end{solution}

\begin{solution}{exo:b3:inorganic-materials:11}
$\dfrac{10n^2 + 2}{(10n^3 + 15n^2 + 11n + 3)/3} = \dfrac{30n^2 + 6}{10n^3 + 15n^2 + 11n + 3} \to \dfrac{30n^2}{10n^3} = \dfrac3n$. Pour $n = 8$~: valeur exacte
$642/2057 = 0.31$, contre $3/8 = 0.375$~; l'approximation surestime la fraction des petits
agrégats, dont les couches internes ne sont pas négligeables.
\end{solution}

\begin{solution}{exo:b3:inorganic-materials:12}
$2E = 3V$, $2E = 5p + 6h + 7s$, $F = p + h + s$. Euler~: $V - E + F = 2$, avec $V = 2E/3$~: $F - E/3 = 2$, donc
$p + h + s - (5p + 6h + 7s)/6 = 2$, c'est-à-dire $(p - s)/6 = 2$~: $p - s = 12$. Chaque heptagone exige
un pentagone supplémentaire.
\end{solution}

\begin{solution}{pb:b3:inorganic-materials:1}
\textbf{1.} $\mathrm{Na_{12}Al_{12}Si_{12}O_{48}}$~: $12 \times 23.0 + 12 \times 27.0 + 12 \times 28.1 + 48 \times 16.0 = \qty{1705.2}{g/mol}$.

\textbf{2.} $1705.2 + 27 \times 18.0 = \qty{2191.2}{g/mol}$.

\textbf{3.} Si/Al $= 12/12 = 1$.

\textbf{4.} Pas de liaisons Al--O--Al~; avec Si/Al $= 1$, chaque Si a quatre voisins Al et
chaque Al quatre voisins Si~: alternance stricte.

\textbf{5.} $-12$, une par aluminium.

\textbf{6.} $486.0/2191.2 = 22.2$\,\%.

\textbf{7.} $\ce{2Na+}(\text{zéolithe}) + \ce{Ca^2+}(\text{aq}) \rightleftharpoons \ce{Ca^2+}(\text{zéolithe}) + \ce{2Na+}(\text{aq})$.

\textbf{8.} Six (douze charges).

\textbf{9.} $6/\qty{1705.2}{g/mol} = \qty{3.52}{mmol/g}$.

\textbf{10.} $6/\qty{2191.2}{g/mol} = \qty{2.74}{mmol/g}$.

\textbf{11.} $\qty{3.52}{mmol/g} \times \qty{100.1}{g/mol} = \qty{352}{mg}$ de $\ce{CaCO3}$ par gramme.

\textbf{12.} $\qty{2.0}{mmol/L} \times \qty{15}{L} = \qty{30}{mmol}$.

\textbf{13.} $\qty{30}{mmol}/\qty{2.74}{mmol/g} = \qty{11}{g}$.

\textbf{14.} $\qty{10.9}{g}/0.60 = \qty{18}{g}$.

\textbf{15.} $\qty{60}{mmol}$ de $\ce{Na+}$, $\qty{0.060}{mol} \times \qty{23.0}{g/mol} = \qty{1.4}{g}$.

\textbf{16.} Les phosphates des eaux usées nourrissent les algues des
rivières et des lacs~; leur prolifération puis leur décomposition
consomment l'oxygène dissous (eutrophisation).

\textbf{17.} L'ion $\ce{Mg^2+}$, plus petit, retient plus fortement ses molécules d'eau~; il doit s'en défaire
pour franchir la fenêtre et atteindre les sites, ce qui est lent aux
températures de lavage.

\textbf{18.} L'ion nu mesure environ \qty{0.20}{nm} de diamètre, la moitié de la fenêtre.

\textbf{19.} L'ion hydraté est plus gros que la fenêtre~: il doit perdre une
partie de sa couche d'hydratation pour passer, et les oxygènes de la
charpente prennent la place de l'eau.

\textbf{20.} La charpente est chargée négativement et n'échange que des
cations~; un gros anion à longue chaîne est repoussé et ne peut pas
franchir la fenêtre, si bien qu'il reste dans l'eau, où il fait son
travail.

\textbf{21.} Déshydratée, la \omterm{def:b3:inorganic-materials:zeolite}{zéolithe} fixe fortement l'eau dans ses
cages~; seules les molécules plus petites que la fenêtre entrent, si bien
qu'elle sépare les molécules selon leur taille~: un tamis à l'échelle
moléculaire.

\textbf{22.} Les ions $\ce{K+}$, plus gros, se logent dans les fenêtres et les rétrécissent~: seules des
molécules plus petites, comme l'eau, sont admises.

\textbf{23.} L'isomère \textit{para}, élancé et linéaire, s'ajuste aux
canaux et diffuse vite~; les isomères \textit{ortho} et \textit{meta}, plus
larges, diffusent lentement, restent et s'isomérisent.

\textbf{24.} La \omterm{def:b3:inorganic-materials:zeolite}{zéolithe} A anhydre peut fixer, en théorie, \qty{3.52}{mmol} de $\ce{Ca^2+}$ par gramme.
\end{solution}
